% GENERATED FILE. Edit canonical lessons, then run npm run generate:books.
% canonical-source-hash: fnv1a64:2b1c65799afc0c51
% canonical-lessons: ES-C507-votar, ES-C507-impuesto, ES-C507-poblacion, ES-C507-habitante, ES-C507-debate

\chapter{La sociedad}
\label{ch:es-la-sociedad}
\hlchaptermodality{507}

\begin{chapteropening}
\textbf{By the end of this chapter:} \emph{I can read about an election, taxes, the population, the inhabitants and a debate.}

Complete the last lesson of chapter 507: I can read about an election, taxes, the population, the inhabitants and a debate.
\end{chapteropening}

\section[votar]{votar — to vote}

\label{lesson:ES-C507-votar}



\begin{quote}
Before the new one: say the Spanish for to congratulate, then the Spanish for to hug.
\end{quote}

\subsection*{You'll want to know: votar}
\textbf{votar} — \textquotedblleft{}to vote\textquotedblright{}.

\begin{quote}
Mañana votamos en el colegio del barrio.
\end{quote}

\begin{quote}
\emph{Tomorrow we vote at the local school.}
\end{quote}

\begin{cousinweb}
\textbf{votar} is built on \emph{voto}, \textquotedblleft{}a vote, a vow\textquotedblright{}, from Latin \emph{votum}, \textquotedblleft{}a vow\textquotedblright{} — a learned word. Worn down in speech, its plural \emph{vota} became \emph{boda}, \textquotedblleft{}a wedding\textquotedblright{}.
\end{cousinweb}

\begin{grammarlens}[title={What you've built}]
One more word for reading notices, forms and messages.
\end{grammarlens}

\subsection*{Your turn}
\begin{itemize}
\raggedright
  \item \emph{Say it:} \emph{votar}
  \item \emph{Say it:} \emph{votar}, once more
  \item \emph{Say it:} a sentence of your own with \emph{votar}
  \item \emph{Recall:} say the Spanish for to congratulate, then the Spanish for to hug, then say \emph{votar} again
\end{itemize}

\begin{tcolorbox}[breakable,colback=teal!4,colframe=teal!35!black,title={Before you move on}]
What does votar mean? (\textquotedblleft{}To vote\textquotedblright{}.) Say it once more.
\end{tcolorbox}
\section[el impuesto]{el impuesto — the tax}

\label{lesson:ES-C507-impuesto}



\begin{quote}
Before the new one: say the Spanish for to hug, then the Spanish for to vote.
\end{quote}

\subsection*{You'll want to know: el impuesto}
\textbf{el impuesto} — \textquotedblleft{}the tax\textquotedblright{}.

\begin{quote}
El precio incluye los impuestos.
\end{quote}

\begin{quote}
\emph{The price includes taxes.}
\end{quote}

\begin{cousinweb}
\textbf{impuesto} is the participle of \emph{imponer}, \textquotedblleft{}to impose\textquotedblright{}, Latin \emph{imponere}, \textquotedblleft{}to place upon\textquotedblright{}: a tax is what is laid on you. \emph{Puesto}, \textquotedblleft{}put\textquotedblright{}, is the inherited participle of \emph{poner}.
\end{cousinweb}

\begin{grammarlens}[title={What you've built}]
One more word for reading notices, forms and messages.
\end{grammarlens}

\subsection*{Your turn}
\begin{itemize}
\raggedright
  \item \emph{Say it:} \emph{el impuesto}
  \item \emph{Say it:} \emph{el impuesto}, once more
  \item \emph{Say it:} a sentence of your own with \emph{el impuesto}
  \item \emph{Recall:} say the Spanish for to hug, then the Spanish for to vote, then say \emph{el impuesto} again
\end{itemize}

\begin{tcolorbox}[breakable,colback=teal!4,colframe=teal!35!black,title={Before you move on}]
What does el impuesto mean? (\textquotedblleft{}The tax\textquotedblright{}.) Say it once more.
\end{tcolorbox}
\section[la población]{la población — the population}

\label{lesson:ES-C507-poblacion}



\begin{quote}
Before the new one: say the Spanish for to vote, then the Spanish for the tax.
\end{quote}

\subsection*{You'll want to know: la población}
\textbf{la población} — \textquotedblleft{}the population\textquotedblright{}.

\begin{quote}
La población del pueblo es pequeña.
\end{quote}

\begin{quote}
\emph{The village's population is small.}
\end{quote}

\begin{cousinweb}
\textbf{población} $\leftarrow$ Late Latin \emph{populatio} (EtymDB 2.1), from \emph{populus}, \textquotedblleft{}a people\textquotedblright{} — the word that also became \emph{pueblo}, \textquotedblleft{}a town, a people\textquotedblright{}.
\end{cousinweb}

\begin{grammarlens}[title={What you've built}]
One more word for reading notices, forms and messages.
\end{grammarlens}

\subsection*{Your turn}
\begin{itemize}
\raggedright
  \item \emph{Say it:} \emph{la población}
  \item \emph{Say it:} \emph{la población}, once more
  \item \emph{Say it:} a sentence of your own with \emph{la población}
  \item \emph{Recall:} say the Spanish for to vote, then the Spanish for the tax, then say \emph{la población} again
\end{itemize}

\begin{tcolorbox}[breakable,colback=teal!4,colframe=teal!35!black,title={Before you move on}]
What does la población mean? (\textquotedblleft{}The population\textquotedblright{}.) Say it once more.
\end{tcolorbox}
\section[el habitante]{el habitante — the inhabitant}

\label{lesson:ES-C507-habitante}



\begin{quote}
Before the new one: say the Spanish for the tax, then the Spanish for the population.
\end{quote}

\subsection*{You'll want to know: el habitante}
\textbf{el habitante} — \textquotedblleft{}the inhabitant\textquotedblright{}.

\begin{quote}
La ciudad tiene un millón de habitantes.
\end{quote}

\begin{quote}
\emph{The city has a million inhabitants.}
\end{quote}

\begin{cousinweb}
\textbf{habitante} $\leftarrow$ Latin \emph{habitans}, \emph{habitantis}, \textquotedblleft{}dwelling\textquotedblright{}, the participle of \emph{habitare}, \textquotedblleft{}to dwell\textquotedblright{} — the \emph{habitare} inside \emph{habitación}. A learned word: the \emph{-t-} between vowels stayed \emph{t}.
\end{cousinweb}

\begin{grammarlens}[title={What you've built}]
One more word for reading notices, forms and messages.
\end{grammarlens}

\subsection*{Your turn}
\begin{itemize}
\raggedright
  \item \emph{Say it:} \emph{el habitante}
  \item \emph{Say it:} \emph{el habitante}, once more
  \item \emph{Say it:} a sentence of your own with \emph{el habitante}
  \item \emph{Recall:} say the Spanish for the tax, then the Spanish for the population, then say \emph{el habitante} again
\end{itemize}

\begin{tcolorbox}[breakable,colback=teal!4,colframe=teal!35!black,title={Before you move on}]
What does el habitante mean? (\textquotedblleft{}The inhabitant\textquotedblright{}.) Say it once more.
\end{tcolorbox}
\section[el debate]{el debate — the debate}

\label{lesson:ES-C507-debate}



\begin{quote}
Before the new one: say the Spanish for the population, then the Spanish for the inhabitant.
\end{quote}

\subsection*{You'll want to know: el debate}
\textbf{el debate} — \textquotedblleft{}the debate\textquotedblright{}.

\begin{quote}
Vemos el debate en la tele.
\end{quote}

\begin{quote}
\emph{We watch the debate on TV.}
\end{quote}

\begin{cousinweb}
\textbf{debate} is the noun of \emph{debatir}, \textquotedblleft{}to argue\textquotedblright{}, built from \emph{de-} and \emph{batir}, \textquotedblleft{}to beat\textquotedblright{}, Latin \emph{battuere}. French \emph{débattre} is built the same way.
\end{cousinweb}

\begin{grammarlens}[title={What you've built}]
One more word for reading notices, forms and messages.
\end{grammarlens}

\subsection*{Your turn}
\begin{itemize}
\raggedright
  \item \emph{Say it:} \emph{el debate}
  \item \emph{Say it:} \emph{el debate}, once more
  \item \emph{Say it:} a sentence of your own with \emph{el debate}
  \item \emph{Recall:} say the Spanish for the population, then the Spanish for the inhabitant, then say \emph{el debate} again
\end{itemize}

\begin{tcolorbox}[breakable,colback=teal!4,colframe=teal!35!black,title={Before you move on}]
What does el debate mean? (\textquotedblleft{}The debate\textquotedblright{}.) Say it once more.
\end{tcolorbox}
